\chapter{Scientific Theories: Revision as Relay Replacement}
\label{ch:science}

\begin{quotation}
\noindent\textit{Synopsis.} Scientific theories are treated here as finite
renderings of an open empirical world. The chapter argues that a theory
inevitably encounters observations it did not anticipate and therefore
cannot remain sufficient indefinitely; scientific progress is reframed as
relay replacement rather than local correction, with a new theory replacing
an old relay once accumulated insufficiency becomes too great to ignore.
Unlike most domains treated in the monograph, science actively depends on
discovering its own insufficiency, making it a particularly revealing test
of the framework.
\end{quotation}

\section{A Theory as a Finite Substrate for an Open Observation Space}

\begin{proposition}[Theories Are Renderings of an Open Observation Space]
\label{prop:theory-rendering}
A scientific theory is a finite rendering of a potentially unbounded
space of future observations.
\end{proposition}

\begin{proofsketch}
The observation space relevant to a theory is not fixed at the theory's
formulation; new instruments, new regimes, and new experimental designs
continually generate previously unavailable observations. A theory,
being a finite statement, cannot represent this space exhaustively, and
\cref{def:rendering}'s projection structure therefore applies.
\end{proofsketch}

\begin{proposition}[Revision Necessity]
\label{prop:revision-necessity}
Every finite theory admits possible future observations that would
require its revision.
\end{proposition}

\begin{proofsketch}
This is \cref{prop:task-unavailability} applied to
\cref{prop:theory-rendering}'s open observation space: a future
observation $\tau^*$ can always be constructed that the theory, as a
finite rendering, was not built to accommodate. This is essentially the
falsifiability criterion, reframed within the admissibility apparatus of
Chapter~\ref{ch:admissibility} rather than derived independently.
\end{proofsketch}

\section{Revision, Not Augmentation}

Unlike chess (\cref{ch:chess}), where insufficiency is resolved by
augmenting a closed hidden-state set, and unlike maps
(\cref{ch:maps}), where insufficiency is resolved by declaring the
discard, a theory's response to \cref{prop:revision-necessity} is
typically wholesale replacement of the theory itself
(\cref{def:relay-revision}, \cref{sec:coord-revision}), not a patch to
the existing substrate. This is the one domain treated in this monograph
where global insufficiency, per \cref{principle:global-insufficiency},
is not merely tolerated but functions as the intended mechanism of
progress: a theory's eventual insufficiency is what licenses its
replacement by a better one.

\begin{example}
A theory calibrated against observations within a certain regime --
low relative velocities, say -- can be sufficient
(\cref{def:sufficiency}) for every task posed within that regime while
being straightforwardly insufficient for tasks posed in a regime its
observational basis never covered. The historical replacement of such a
theory by one validated across a wider observational range is a
canonical instance of \cref{def:relay-revision}: the newer theory is not
derived from the older one by augmentation, in the sense
\cref{prop:closed-sufficiency} would require, since the older theory's
observational basis was never closed to begin with; it is a new
production built from an enlarged evidential base.
\end{example}

\section{A Note on Adjacent Debates}

This chapter's claim is narrower than it might appear next to some
well-known positions in philosophy of science, and it is worth marking
the boundary rather than letting the two blur together.
\cref{prop:revision-necessity} claims only that a finite theory cannot
anticipate every future observation, which is a direct instance of
\cref{prop:task-unavailability} and carries no commitment about whether
successive theories are commensurable, whether observation itself is
theory-laden in ways that complicate what counts as a ``new''
observation, or whether scientific change is better described by gradual
accumulation or by more discontinuous conceptual shifts. These are
substantive, contested questions in their own right, and this monograph
does not adjudicate them; \cref{def:relay-revision}'s formal claim --
that a superseded chain is replaced rather than repaired -- is compatible
with more than one position on how sharp or continuous that replacement
is. Chapter~\ref{ch:four-operations}
(\cref{prop:relay-replacement-partial-revision}) draws the same boundary
from the other direction, formally: relay revision corresponds to
Distinguishability Geometry's revision operation only when the underlying
state space is held fixed and the ontology merely expands. The sharper
case, where the state space itself changes, is not left uncharacterized:
\cref{prop:relay-replacement-paradigm-shift} shows it corresponds instead
to Observability Is Restorability's Paradigm Shift Theorem
(\cref{thm:paradigm-shift}), under which the two eras' restorable
distinctions become genuinely incomparable rather than one refining the
other. This monograph still does not claim to have resolved which of
science's historical episodes are which; it claims only that both
candidate cases now have a formal home.

\section{Scientific Revolutions as Relay Replacement}

It is worth drawing the last three sections together into a single
picture, since they were developed at different points in this
monograph and can otherwise read as three separate remarks rather than
one account. A scientific tradition is a substrate chain in the sense of
Chapter~\ref{ch:budget-relay}: successive theoretical and observational
stages, each locally admissible with respect to the evidential standards
and instruments available at the time. \cref{prop:revision-necessity}
guarantees this chain eventually meets an observation it cannot
accommodate. What happens next depends on which of two cases obtains.
If the anomaly can be absorbed by extending the evidential ontology
without disturbing what counts as a state of the system under study --
a more precise instrument within the same conceptual framework --
\cref{prop:relay-replacement-partial-revision} applies, Distinguishability
Geometry's Revision Monotonicity theorem governs the transition, and the
new theory's coding deficit is guaranteed no worse than the old one's:
this is normal science, in something close to the textbook sense,
absorbing anomalies by controlled extension. If instead the anomaly
cannot be absorbed this way -- if accommodating it requires changing
what counts as a state at all -- \cref{prop:relay-replacement-paradigm-shift}
applies instead, and Observability Is Restorability's Paradigm Shift
Theorem characterizes the result: the old and new theories' restorable
distinctions become genuinely incomparable, with real losses alongside
the real gains, rather than the new theory simply subsuming the old one.
Both cases are relay replacement in \cref{def:relay-revision}'s sense;
what distinguishes them is not sharpness of the break in some vague,
qualitative sense but a precise, checkable structural fact --- whether
the state space was held fixed. This monograph does not attempt to sort
particular historical episodes into one case or the other, since doing
so responsibly would require the kind of close historical argument this
book's method does not supply; what it offers instead is the distinction
itself, sharp enough that such sorting could in principle be attempted
by those equipped to attempt it.
